\documentclass[12pt,a4paper]{article}
\usepackage[a4paper,margin=2.1cm]{geometry}
\usepackage{fontspec}
\usepackage{polyglossia}
\setdefaultlanguage{english}
\setotherlanguage{hebrew}
\setmainfont{Times New Roman}
\newfontfamily\hebrewfont[Script=Hebrew,Scale=1.06]{Arial Unicode MS}
\newfontfamily\hebrewfonttt[Script=Hebrew,Scale=1.0]{Arial Unicode MS}
\usepackage{xcolor}
\usepackage{enumitem}
\usepackage{hyperref}
\usepackage{titlesec}
\usepackage{setspace}
\setstretch{1.12}

\hypersetup{
  colorlinks=true,
  linkcolor=blue!40!black,
  urlcolor=blue!40!black
}

\titleformat{\section}{\Large\bfseries}{}{0pt}{}
\titleformat{\subsection}{\normalsize\bfseries}{}{0pt}{}
\setlist[itemize]{itemsep=3pt,topsep=4pt}

\title{\texthebrew{דף מידע לשאלות אחרי המצגת}\\[4pt]\large MHM Pipeline}
\author{Alexander Goldberg}
\date{\texthebrew{מאי 2026}}

\begin{document}
\maketitle

\begin{hebrew}
זהו דף צדדי, לא טקסט דיבור. המטרה שלו היא לתת תשובות קצרות וברורות לשאלות שיכולות לעלות אחרי ההרצאה.
\end{hebrew}

\section{\texthebrew{מה המקור לטענה של כ-95\%?}}
\begin{hebrew}
אם שואלים על המספר הזה, כדאי לענות בזהירות:
\par
הכוונה אינה לכך שכל הנתונים שלי מכסים 95\% מן העולם, אלא להערכה קטלוגית רחבה לגבי היקף פרויקט כתבי היד העבריים.
\par
המקור הוא עמוד ההסבר של הספרייה הלאומית על פרויקט כתיב. שם נאמר שהמכון לתצלומי כתבי יד עבריים אסף עותקי מיקרופילם של קרוב ל-95\% מכתבי היד העבריים הידועים בעולם.
\par
ניסוח טוב בעל פה:
\par
``כשאני אומר כ-95\%, אני מתייחס להערכת הכיסוי של המכון לתצלומי כתבי יד עבריים, כפי שהיא מופיעה בתיאור הרשמי של פרויקט כתיב של הספרייה הלאומית. זה מתאר את היקף התשתית הקטלוגית, לא את מערך המדידה של הפייפליין.''
\par
\textbf{Reference:} National Library of Israel, About the Ktiv Project:
\url{https://www.nli.org.il/en/discover/manuscripts/hebrew-manuscripts/about}
\end{hebrew}

\section{\texthebrew{מה ההבדל בין 50 אלף רשומות לבין 68 רשומות?}}
\begin{hebrew}
יש כאן שני שימושים שונים בנתונים:
\begin{itemize}
  \item בערך 50 אלף רשומות שימשו לניתוח ראשוני של הקורפוס: אילו שדות קיימים, כמה טקסט חופשי יש, ואיפה חסר מידע.
  \item 68 רשומות שימשו כקבוצת בדיקה קבועה. עליהן נמדדים המספרים שאני מציג בהרצאה ובמאמר.
\end{itemize}
\par
ניסוח טוב בעל פה:
\par
``קורפוס המחקר הרחב עזר לי להבין איך נראות הרשומות. אבל כדי למדוד את הפייפליין בצורה עקבית, השתמשתי בקבוצת בדיקה קטנה וקבועה של 68 רשומות, שנבחרו כדי לכלול הרבה סוגים של מצבים קטלוגיים.''
\end{hebrew}

\section{\texthebrew{למה דווקא 68 רשומות?}}
\begin{hebrew}
המספר 68 לא נבחר מראש כמספר עגול. הוא יצא מתהליך בחירה קבוע שאפשר לשחזר.
\par
המטרה הייתה לקבל קבוצת בדיקה קטנה מספיק לבדיקה ידנית, אבל עשירה מספיק כדי לבדוק הרבה מקרים שונים.
\par
התהליך בקצרה:
\begin{itemize}
  \item 8 רשומות נדירות: מקרים שקל לפספס, כמו לאדינו, יוונית, תאריך במאה עברית, או הערת בעלות ארוכה.
  \item 9 רשומות שמוסיפות כיסוי: בכל פעם נבחרה הרשומה שמוסיפה הכי הרבה תופעות שעדיין לא הופיעו.
  \item 43 רשומות טיפוסיות: כדי שקבוצת הבדיקה לא תהיה רק אוסף של חריגים.
  \item 8 רשומות משלימות: מקרים מיוחדים שלא הופיעו במדגם הראשי.
\end{itemize}
\par
המקור הטכני הוא \texttt{data/tsvs/test\_subset.tsv}, שנוצר על ידי \texttt{scripts/build\_test\_subset.py}. המניפסט של הקובץ מתעד את ארבעת השלבים ואת חתימת ה-SHA של הקורפוס.
\end{hebrew}

\section{\texthebrew{מה זה "סיגנל קשה"?}}
\begin{hebrew}
סיגנל הוא תופעה קטלוגית שרצינו שקבוצת הבדיקה תכיל לפחות פעם אחת.
\par
סיגנל קשה הוא תופעה נדירה שקל לפספס בדגימה רגילה.
\par
דוגמאות:
\begin{itemize}
  \item תאריך לפני 1582, שמצריך טיפול בלוח יוליאני.
  \item שפה נדירה יחסית בקורפוס, כמו יוונית או לאדינו.
  \item שדה 880, שמייצג צורת כתיבה חלופית.
  \item הערת בעלות ארוכה בשדה 561.
  \item כותרת כללית כמו ``קובץ'', שקשה לפרש באופן אוטומטי.
\end{itemize}
\par
הסיבה שהכרחנו את המקרים האלה להיכנס היא שאם לא עושים את זה, דגימה רגילה עלולה להיראות טובה מדי, כי היא לא בודקת את המקרים הקשים באמת.
\end{hebrew}

\section{\texthebrew{מה הם שדות ה-MARC שהוזכרו בהרצאה?}}
\begin{hebrew}
\begin{itemize}
  \item \textbf{MARC 500}: הערות כלליות. כאן מופיעים הרבה פרטים חופשיים, למשל קולופון, תיאור פיזי, או מידע שהקטלוג לא ידע לשים בשדה מובנה.
  \item \textbf{MARC 505}: תוכן כתב היד. למשל אילו חיבורים נמצאים בו, באילו דפים, ולעיתים מי המחבר.
  \item \textbf{MARC 561}: בעלות ותולדות מעבר. למשל מי החזיק בכתב היד, באיזה אוסף הוא היה, או מאיזה מקום הגיע.
  \item \textbf{MARC 655}: ז'אנר או סוג החיבור. למשל תפילה, קבלה, פירוש, שירה.
  \item \textbf{MARC 100/700}: אנשים ותפקידים. 100 הוא בדרך כלל אדם מרכזי ברשומה; 700 הוא אדם נוסף שקשור לרשומה. שלב הסמכות צורך גם 110/111/710/711 (ארגונים וכנסים), וכן 600/610 (אנשים וארגונים כנושאים), 651/752 (מקומות), ו-800/810/811 (כניסות סדרה) --- שלוש-עשרה תגיות בסך הכול.
  \item \textbf{MARC 008}: שדה מקודד וקצר שכולל בין השאר מידע על תאריך ושפה.
\end{itemize}
\end{hebrew}

\section{\texthebrew{אילו סוגי מידע אנחנו מחלצים מן ה-MARC?}}
\begin{hebrew}
בהרצאה מספיק לומר שאנחנו מחלצים שמות אנשים, תאריכים ותוכן של כתבי יד. אם שואלים בפירוט, אפשר להרחיב:
\begin{itemize}
  \item כותרות ושמות חיבורים.
  \item אנשים ותפקידים, למשל מחבר, מעתיק, בעלים, מתרגם או מפרש.
  \item תאריכים ושפות.
  \item תוכן כתב היד: אילו חיבורים מופיעים בו, באילו דפים, ולעיתים מי המחברים שלהם.
  \item בעלות, אוספים ותולדות מעבר של כתב היד.
  \item נושאים, ז'אנרים, קישורים ומזהים חיצוניים.
  \item לעיתים גם מאפיינים פיזיים, כמו חומר, מידות או עיטור.
\end{itemize}
\par
חשוב להדגיש: לא כל רשומה מכילה את כל הסוגים האלה, ולא כל מידע כזה עובר אוטומטית לוויקינתונים. הפייפליין קודם מחלץ ובודק, ורק אחר כך מכין מידע שאפשר לעבור עליו.
\end{hebrew}

\section{\texthebrew{מה זה NER ומה זה BIO?}}
\begin{hebrew}
NER הוא קיצור של \textenglish{Named Entity Recognition}, ובעברית: זיהוי ישויות בשם.
\par
הרעיון הוא לתת למודל טקסט חופשי ולבקש ממנו לסמן בתוכו קטעים בעלי משמעות: אדם, בעלים, חיבור, תאריך, אוסף או סימון דף.
\par
כדי שהמודל ידע איפה כל ישות מתחילה ואיפה היא נגמרת, משתמשים בשיטת סימון שנקראת BIO:
\begin{itemize}
  \item \textbf{B} מסמן את תחילת הישות.
  \item \textbf{I} מסמן המשך של אותה ישות.
  \item \textbf{O} מסמן מילה שאינה חלק משום ישות.
\end{itemize}
\par
לדוגמה, בשם ``רבי משה בן מימון'', אפשר לסמן:
\par
\textenglish{B-PERSON | I-PERSON | I-PERSON | I-PERSON}
\par
כלומר, המודל לא רק אומר שיש כאן אדם, אלא גם מסמן את גבולות הביטוי.
\end{hebrew}

\section{\texthebrew{מה היו שקפי הדוגמאות הטכניים שהועברו ל-Q\&A?}}
\begin{hebrew}
הורדנו מן המצגת את פירוט הדוגמאות כדי לשמור על קצב של עשר דקות, אבל אפשר להשתמש בהן בתשובות:
\begin{itemize}
  \item \textbf{Person NER}: מקבל משפט או קטע הערה, ומסמן שמות אנשים ותפקידים.
  \item \textbf{Ownership NER}: מקבל טקסט שעוסק בבעלות ותולדות מעבר, ומסמן בעלים, תאריכים ואוספים.
  \item \textbf{Contents NER}: מקבל תיאור תוכן של כתב יד, ומסמן שם חיבור, מחבר חיבור וסימון דף.
  \item \textbf{Genre classifier}: מקבל כותרת והקשר מקומי מתוך הרשומה, ומציע ז'אנר כאשר שדה הז'אנר חסר או לא מספיק ברור.
\end{itemize}
\par
ניסוח טוב בעל פה:
\par
``במצגת לא נכנסתי לסימון הטוקנים עצמו, אבל טכנית המודלים עובדים כך: חלקם מסמנים מילים בתוך משפט, וחלקם מסווגים משפט או רשומה. זאת הסיבה שאנחנו משתמשים בכמה מודלים ממוקדים ולא במודל אחד כללי.''
\end{hebrew}

\section{\texthebrew{מאיפה מגיע השיפור בכיסוי?}}
\begin{hebrew}
חשוב לא לייחס את כל השיפור למודלי NER. כל שכבה תורמת משהו אחר:
\begin{itemize}
  \item \textbf{המרה ישירה של MARC}: זה קו הבסיס. הוא מייצר בערך 8 טענות לכל כתב יד, ושומר רק מה שכבר קיים בצורה מפורשת ומובנית.
  \item \textbf{חוקים בקוד}: משפרים בעיקר תאריכים. בכיסוי תאריך יצירה, העלייה היא מ-22\% ל-97\%.
  \item \textbf{NER}: מוסיף מידע מתוך טקסט חופשי. למשל, מודל האנשים מפיק 59 מופעי אדם-תפקיד ב-36 מתוך 68 הרשומות.
  \item \textbf{מסווגים}: מסווג הז'אנר מרחיב את כיסוי הז'אנרים מ-28 מתוך 68 רשומות ל-39 מתוך 68 רשומות, כלומר בערך מ-41\% ל-57\%. מסווג MARC 500 שולח משפטי בעלות למודל שמחלץ בעלים, תאריכים ואוספים.
  \item \textbf{מאגרי סמכות (Mazal, VIAF, Wikidata, KIMA)}: לא מוסיפים תוכן חדש על כתב היד, אלא מעשירים זהויות אחרי שכבר זיהינו אותן. VIAF מספק מזהים חיצוניים מספריות לאומיות (GND, LCCN, ISNI, BnF) דרך ``cluster harvesting'' --- התאמה אחת מניבה עד ארבעה מזהים נוספים. Wikidata מספק QID, וכאשר אין שנות לידה/פטירה ממקור אחר --- הצינור משלים אותן ישירות מ-Wikidata (P569/P570) כדי שמשמר התאריכים בשלב 3 יוכל לפעול גם על מועמדים שזוהו רק דרך SPARQL.
\end{itemize}
\par
כשאומרים ``מ-x ל-y'', ה-x הוא לא כל המידע שקיים בקטלוג. הוא קו הבסיס: מה היה אפשר להוציא בהמרה ישירה או לפני השכבה הספציפית שאנחנו בודקים. ה-y הוא התוצאה אחרי הפייפליין המלא, או אחרי אותה שכבה.
\par
אם שואלים ספציפית על תאריכים:
\par
העלייה בכיסוי תאריכים מגיעה בעיקר מחוקים ומפרשים ייעודיים, לא מ-NER. הם קוראים תאריך מובנה כאשר הוא קיים, וגם מזהים ניסוחים כמו ``מאה ט"ז'' או תאריכים חלקיים.
\par
אם שואלים ספציפית על ז'אנרים:
\par
הבסיס הוא MARC 655: ב-28 מתוך 68 רשומות יש ז'אנר מקורי שאפשר למפות. כאשר השדה חסר או לא מספיק, נכנס מסווג הז'אנרים, שמשתמש בכותרת ובהקשר מתוך שדה 500. אחרי השכבה הזאת יש 39 מתוך 68 רשומות עם טענת ז'אנר סופית.
\end{hebrew}

\section{\texthebrew{למה לא להסתפק בחוקים?}}
\begin{hebrew}
חוקים טובים כשיש תבנית יציבה. למשל תאריכים, שדות מקודדים, או מיפוי ידוע בין מונח קטלוגי למזהה.
\par
אבל בשדות חופשיים אין תבנית אחת. אותו אדם או אותו בעלים יכולים להופיע בעשרות ניסוחים שונים, בעברית, בארמית, באנגלית, בקיצורים או במשפטים ארוכים.
\par
לכן חוקים בלבד היו מפספסים הרבה מידע, או דורשים רשימת חריגים אינסופית.
\par
הפתרון הוא שיטה משולבת: חוקים איפה שיש מבנה, מודלים איפה שיש טקסט חופשי, ומאגרי סמכות איפה שצריך לזהות מי האדם או המקום.
\end{hebrew}

\section{\texthebrew{למה לא להשתמש רק במודל שפה גדול?}}
\begin{hebrew}
מודל שפה גדול יכול לעזור בהבנה כללית, אבל הוא פחות מתאים כבסיס יחיד לפייפליין כזה.
\par
יש כאן שלוש סיבות:
\begin{itemize}
  \item צריך תוצאות שחוזרות על עצמן, לא ניסוח חדש בכל ריצה.
  \item צריך לדעת מאיזה שדה ומאיזה משפט כל טענה הגיעה.
  \item צריך להימנע מהשלמות יצירתיות, כי טעות קטנה יכולה להפוך לטענה היסטורית שגויה.
\end{itemize}
\par
לכן המודלים כאן ממוקדים: הם לא מנסים לכתוב תשובה כללית, אלא לסמן טוקנים, לסווג משפטים, או להשלים ז'אנר במקרים מוגדרים.
\end{hebrew}

\section{\texthebrew{מה זה סוכן ההערכה (eval-agent), ואיך עובד Gemini-כשופט?}}
\begin{hebrew}
סוכן ההערכה הוא רכיב נפרד מן הצינור עצמו. הוא קורא את ה-JSON שיצא משלב החילוץ, וקלט נוסף --- רשומות ה-MARC המקוריות --- ומשווה ביניהם.
\par
לכל ישות שחולצה, הסוכן בונה פרומפט שמציג למודל \emph{Gemini} שלושה דברים: רשומת ה-MARC כפי שהיא, הישות שחולצה, והקריטריונים לפי סוג המסווג (NER לאנשים, NER לבעלות, מסווג ז'אנר וכדומה). Gemini מחזיר אחת מארבע אפשרויות: ``נראה נכון'', ``שגוי'', ``נכון חלקית'' או ``לא ניתן לקבוע''.
\par
הסוכן רץ במקביל על מספר מסווגים, שומר כל פסק דין ב-cache עם תיוג של הרשומה, המסווג, וגרסת הפרומפט. בסוף הריצה מתקבל דו"ח דיוק לכל מסווג --- כמה הוא צדק, כמה טעה, וכמה נכון חלקית.
\par
חידוש שהוספנו לאחרונה הוא עמודת ``מידע חדש'': כשהסוכן מאשר טענה, אנחנו בודקים אוטומטית אם הערך מופיע כבר בשדות ה-MARC המקודדים. אם לא --- הוא מסומן כמידע שהפייפליין הוסיף לקטלוג, ולא שיחזר ממנו.
\par
חידוש נוסף: \textbf{תיקון אוטומטי מוצע (Auto-fix)}. כאשר הסוכן מגיע למסקנה שישות NER \emph{נכונה חלקית} --- למשל שם אדם שמפסיד פטרונים שנוכח בבירור בשדה MARC --- הוא יכול להציע טקסט מתוקן. הצעה כזאת, ברמת ביטחון \textenglish{high} בלבד, מוצגת כלחצן \textbf{Auto-fix} בצד השורה הרלוונטית בטבלת הישויות ובמגירת הפרטים. לחיצה עליו שולחת PATCH שמעדכן את הטקסט מבלי לאשר את הישות --- ההחלטה לאשר נשארת בידי האדם. מגבלות הבטיחות: הטקסט המוצע חייב להיות שונה מהמקורי, לא ריק, לא יותר מ-512 תווים, ויש צורך בהקשר MARC או בסיגנל עיגון מהסטטוס; ישויות מסווג הז'אנר אינן נכנסות לנתיב הזה כלל.
\end{hebrew}


\section{\texthebrew{האם נכון מתודולוגית להשתמש ב-LLM כדי להעריך את הפלט של מודל אחר?}}
\begin{hebrew}
זאת שאלה לגיטימית, וצריך לענות עליה במדויק: \emph{הסוכן אינו תחליף לסקירה אנושית; הוא שכבת ביקורת בקנה מידה.}
\par
כששואלים את Gemini ``האם הישות שחולצה תואמת את רשומת ה-MARC?'' --- זאת לא משימה יצירתית. זאת משימה של אימות מול עוגן טקסטואלי, שמודלים גדולים מבצעים יחסית טוב. אנחנו לא מבקשים מ-Gemini לדרג איכות, אלא לסמן אי-התאמות בולטות.
\par
שנית, השוואה אנושית. עבור תת-מדגם של 100 פסקי דין השוויתי ידנית את הסיווג של הסוכן מול שיפוט מומחה. הסכמה הכוללת הייתה כ-92\%; רוב הפערים היו מקרי גבול שגם בין שני שופטים אנושיים היו מקבלים סיווג שונה.
\par
לבסוף, הפסקי-דין של הסוכן הם \emph{עוזרים}, לא קובעים: כל טענה שמתפרסמת לויקינתונים עוברת דרך Wikidata Studio --- שם החוקר רואה גם את הראיות וגם את פסק הדין, ומאשר ידנית.
\end{hebrew}


\section{\texthebrew{מה קורה כש-Gemini טועה בפסק הדין?}}
\begin{hebrew}
זה קורה, וזה לא מאיים על המתודולוגיה. שלוש שכבות הגנה:
\par
ראשית, השופט אינו הסמכות הסופית. הסוכן יוצר \emph{המלצה} לחוקר; ההחלטה הסופית על פרסום מתקבלת ב-Wikidata Studio, עם כל הראיות פרושות.
\par
שנית, יש מתג ``re-check (ignore cache)'' שמאפשר להריץ את אותה ישות מחדש דרך Gemini, ולקבל פסק דין חדש. אם פסק הדין הראשון נראה חשוד, אפשר לבדוק שוב.
\par
שלישית, פסקי הדין שמורים עם גרסת הפרומפט וגרסת המודל. אם בעתיד נגלה שגרסת Gemini מסוימת מטה את פסק הדין באופן שיטתי, אפשר להריץ מחדש רק את הפריטים שהסתמכו על אותה גרסה.
\par
חשוב להדגיש: מה שאנחנו מודדים זה \emph{כיסוי} ו\emph{דיוק בקנה מידה}. עבור טענה בודדת, רק החוקר האנושי יכול לאשר.
\end{hebrew}


\section{\texthebrew{איך סוכן הבינה יכול להעריך גם את התאמות הסמכות, ולא רק את החילוץ?}}
\begin{hebrew}
בתחילה הסוכן העריך רק את שלב החילוץ --- את שלושת מודלי ה-NER ואת מסווג הז'אנר. לאחרונה הוספנו מעריך \emph{סמכות} (authority) שבוחן את שלב 3.
\par
ההבדל בקלט: בעוד מעריך החילוץ קורא את \texthebrew{ner\_results.json}, מעריך הסמכות קורא את \texthebrew{authority\_enriched.json} --- הרשומה שכבר נושאת את שדה \texthebrew{marc\_authority\_matches}. עבור כל התאמה שהפייפליין ביצע --- Mazal, VIAF, ויקינתונים או KIMA --- הסוכן שואל: בהינתן הקשר ה-MARC של כתב היד, האם זו אכן רשומת הסמכות הנכונה לשם הזה? המשימה אינה ``לדרג איכות'' אלא לסמן אי-התאמות --- אדם שמת מאות שנים לפני הכתיבה, מקום שלא תואם, או מזהה סמכות שצורף לשם שאינו תואם.
\par
מבחינת ממשק, כפתור ``אימות עם סוכן הבינה'' מופיע כעת גם בלוח הסמכות, ולא רק בלוח החילוץ. אותו דו"ח דיוק, אותם ארבעה פסקי דין, אותו cache --- רק מקור קלט אחר.
\end{hebrew}


\section{\texthebrew{מה זה ``אישור משותף''? אם אני מאשר במסך הבינה, זה משנה את העריכות שלי בקטלוג?}}
\begin{hebrew}
``אישור משותף'' פירושו שיש מצב אישור \emph{אחד} בלבד, נגיש משלושה מסכים שונים: עורך הישויות, עורך הסמכות, ומסך הביקורת של הסוכן. כולם קוראים וכותבים לאותו קובץ סיידקאר, \texthebrew{approvals.json}, שיושב בתיקיית הפלט.
\par
המפתח הוא קנוני ובלתי תלוי במסך: \texthebrew{control-number} של הרשומה, הקבוצה (חילוץ או סמכות), התפקיד או הסוג, וטקסט הישות לאחר נרמול. כך שאם אתה מאשר בעורך הסמכות, ואותה ישות מופיעה גם במסך הסוכן --- שני המסכים מצביעים על אותה רשומה בקובץ.
\par
הסנכרון הוא בזמן אמת. \texthebrew{QFileSystemWatcher} מאזין לקובץ, וכשמסך אחד כותב --- המסכים האחרים שכבר פתוחים טוענים מחדש מיידית. אתה מאשר פעם אחת ורואה את זה בכל מקום. אין מצב שבו ``אישרתי כאן'' ו``זה לא נשמר שם''. וכן --- זה משפיע על מה שיתפרסם: בשמירה, קבצי התוצאה עדיין מסוננים לטענות המאושרות בלבד, בדיוק כמו קודם. חוזה הפרסום לוויקינתונים לא השתנה; רק מקורות האישור התאחדו.
\end{hebrew}


\section{\texthebrew{האם האישור של הבינה הוא אותו אישור כמו שלי?}}
\begin{hebrew}
לא. הסוכן \emph{מציע}; האדם \emph{מכריע}. פסק הדין של הסוכן (``נראה נכון'' וכדומה) הוא המלצה שעוזרת לחוקר לתעדף --- הוא צובע את השורות ומסמן את החשודות, אבל הוא אינו מסמן אותן כ``מאושר''.
\par
האישור עצמו הוא פעולה אנושית: לחיצה על ``אשר''/``דחה'' בכל אחד משלושת המסכים. מצב האישור הזה הוא דבר אחד משותף --- בין אם החוקר הגיע אליו אחרי שקרא את המלצת הסוכן, ובין אם הגיע אליו ישירות בעורך. הסוכן יכול לחסוך לחוקר זמן (``תסתכל קודם על אלה שסומנו כשגויים''), אבל אף טענה אינה מתפרסמת רק משום שהסוכן אישר אותה. ההכרעה הסופית, וחתימת ה``מאושר'', תמיד אנושית.
\par
אותו עיקרון חל גם על לחצן ה-Auto-fix: גם כשהסוכן מציע תיקון ברמת ביטחון גבוהה, הלחיצה על האוטו-פיקס רק מעדכנת את הטקסט --- היא אינה מסמנת את הישות כ``מאושרת''. השדה \texttt{approved} נשאר ריק עד שהחוקר מאשר ידנית.
\end{hebrew}

\section{\texthebrew{מה בדיוק הסוכן יכול לתקן ב-Auto-fix, ומה הוא לא יכול לתקן?}}
\begin{hebrew}
ה-Auto-fix פועל רק על שלושה מסווגי NER: אנשים, פרובננס, ותוכן. מסווג הז'אנר אינו נכנס לנתיב הזה כלל, מפני שז'אנר הוא תווית קטגורית ולא טקסט ישות.
\par
מה שהסוכן יכול לתקן: שם אדם שמפסיד פטרונים הנראים בשדה 100/700 של ה-MARC (``יוסף'' במקום ``יוסף בן יעקב''); שם בעלים שנכתב בצורה מקוצרת; טווח דפים שחולץ חלקית.
\par
מה שהסוכן לא יכול לתקן: טענות שאין להן הקשר MARC שיאמת את התיקון (הסוכן יפיל אז את ה-\textenglish{suggested\_fix}); ישויות שרמת הביטחון של התיקון אינה \textenglish{high}; טקסט מוצע שזהה לטקסט המקורי; טקסט ארוך מ-512 תווים.
\par
ניסוח טוב בעל פה:
\par
``הסוכן יכול לתקן ישויות NER שהיה להן עוגן ב-MARC --- שם שהתמחיש שחולץ ולא שם שהמציא. הלחצן מופיע רק כשהביטחון גבוה וכשאנחנו יכולים להצביע על שדה MARC ספציפי שתמך בתיקון.''
\end{hebrew}


\section{\texthebrew{איך עמודת ``מידע חדש'' מחליטה מה באמת חדש?}}
\begin{hebrew}
זאת בדיקה דטרמיניסטית, לא הסתברותית. עבור כל טענה שהסוכן סימן ``נראה נכון'', הסקריפט בודק האם הערך שחולץ מהטקסט החופשי מופיע כבר באחד מהשדות המובְנים של אותה רשומת MARC --- 100, 110, 700, 710 לאנשים, 651, 752 למקומות, 655 לז'אנר, ועוד.
\par
ההשוואה היא לאחר ניקוי טקסט: הסרת פיסוק, נרמול ניקוד עברי, ובדיקה דו-כיוונית של מחרוזות-משנה (אם השם הוא ``יוסף קארו'' והוא מופיע ב-100 כ``קארו, יוסף בן אפרים'' --- זה נספר כקיים, כי המילים העיקריות חופפות).
\par
אם הערך לא מופיע באף שדה מובְנה --- ✨ ``מידע חדש''. החוקר רואה בעמודה הזאת לא רק שהפייפליין צדק, אלא גם שהוא תרם משהו לקטלוג, ולא רק חזר עליו.
\par
המגבלה הברורה: זאת בדיקה תחבירית, לא סמנטית. אם אותו אדם מופיע ב-MARC בכינוי שונה לחלוטין (``רש"י'' לעומת ``שלמה בן יצחק''), הבדיקה תסמן את החילוץ כחדש גם אם הוא לא באמת חדש. החוקר רואה את שתי הצורות זו לצד זו ויכול להחליט.
\end{hebrew}


\section{\texthebrew{מה העלות והזמן של הרצת הסוכן על אלף רשומות?}}
\begin{hebrew}
המדידות שלנו על מערך הבדיקה (68 רשומות, כ-400 ישויות מחולצות):
\par
זמן: בערך 6 דקות בהרצה רגילה עם 4 פרומפטים מקבילים. הצוואר-בקבוק הוא רחוב ה-API של Google.
\par
עלות API: כ-3 דולר עבור 68 רשומות, לפי מחירון Gemini Flash הציבורי בזמן המדידה. במונחי מחיר-לרשומה: בערך 4 סנט. עבור קורפוס של 1000 רשומות, נצפה לעלות של כ-45 דולר ולזמן של כשעה וחצי.
\par
חידוד חשוב: בזכות ה-cache, ריצה חוזרת על אותו קורפוס (למשל אחרי תיקון פרומפט אצל מסווג בודד) דורשת חלק קטן מן הזמן. רק הישויות שעבורן הפרומפט השתנה ייפנו ל-Gemini.
\par
זה גובה שמתאים לסקירת איכות, לא לפרסום ייצור --- בדיוק כפי שאנחנו ממקמים את הכלי.
\end{hebrew}


\section{\texthebrew{האם תוצאות הסוכן ניתנות לשחזור על ידי חוקרים אחרים?}}
\begin{hebrew}
כן, וזה היה אחד משיקולי העיצוב. שלוש שכבות שחזור:
\par
ראשית, פסקי הדין נשמרים כ-JSONL --- שורה לכל פסק דין, עם שדות מובְנים: מזהה הרשומה, סוג המסווג, גרסת הפרומפט, גרסת המודל, ופסק הדין עצמו. הקובץ הזה ניתן לפרסום עם המאמר.
\par
שנית, הפרומפטים נשמרים בקוד הסוכן כקבצי טקסט סטטיים עם תיוג גרסה. חוקר שמשתחזר את ההערצה מקבל בדיוק את אותו פרומפט.
\par
שלישית, מודל Gemini עצמו נקוב בגרסה. כאשר Google מוציאה גרסה חדשה, אנחנו לא מחליפים אוטומטית --- צריך החלטה מודעת ושחזור של הריצה.
\par
המגבלה היחידה: השופט עצמו (Gemini) אינו פתוח. חוקר שירצה להריץ את אותה הערכה דרך מודל פתוח (למשל Llama-3) יקבל מספרים דומים, אבל לא זהים. זאת מגבלה של שירותי API, לא של המתודולוגיה.
\end{hebrew}


\section{\texthebrew{כמה מהאונטולוגיה באמת מגיעה לוויקינתונים?}}
\begin{hebrew}
זאת שאלה חשובה כי היא בודקת אם הפייפליין באמת ``מקרין'' את הגרף המחקרי, או רק חותך ממנו חתיכה.
\par
המדידה הספציפית היא ברמת מחלקות באונטולוגיה. בקורפוס הבדיקה (68 כתבי יד, 39 מחלקות פעילות):
\begin{itemize}
  \item כ-7.7\% מהמחלקות מקבלות פריט עצמאי בוויקינתונים (P31), כמו כתב יד, אדם, חיבור, מוסד.
  \item כ-79.5\% מקבלות ייצוג ככשרים על טענות קיימות, כלומר מאפייני עזר. דוגמאות: יחידה קודיקולוגית הופכת ל-P2635 (מספר חלקים), קולופון הופך ל-P1684 עם תפקיד, אירוע ייצור הופך ל-P571 + P1071.
  \item כ-12.8\% נשארות בכוונה רק באונטולוגיה ולא עוברות לוויקינתונים: text tradition, transmission witness, paradigm bridge, philological view, evidence chain. אלו מבנים פילולוגיים-פרשניים שמודל הנתונים של ויקינתונים לא בנוי להחזיק.
  \item כיום אין מחלקות לא ממופות בכלל (unknown = 0).
\end{itemize}
\par
הסיכום: כ-87 אחוז מהמחלקות מגיעות לוויקינתונים בצורה כזו או אחרת. ה-13 אחוז שנשארים הם בחירה מודעת, לא חוסר.
\par
חשוב להדגיש שיש שני מסלולי גישה נוספים. מעל הכיסוי הזה, כל פריט בוויקינתונים מקבל מאפיין P2888 שמצביע על מזהה קבוע (w3id.org/mhm) שמפנה לוויקיבייס של הפרויקט; P973 שמוביל לעמוד הוויקיבייס ישירות; ו-P6108 עם מַניפֶסט IIIF Presentation 3.0 שמכיל דף לכל פוליו, טווחים ליחידות קודיקולוגיות, ושכבות הערות לקולופונים ולהתערבויות סופרים. במילים אחרות: מי שמתחיל מפריט בוויקינתונים יכול להגיע בלחיצה אחת אל הגרף המחקרי המלא.
\end{hebrew}


\section{\texthebrew{איך מוודאים שלא חוזרים על הטעויות שגרמו לבעיות מול הקהילה?}}
\begin{hebrew}
זאת שאלה לגיטימית כי במהלך הפיתוח, באפריל 2026, אכן היו תקריות עם הקהילה של ויקינתונים: שגיאות זיהוי, מיזוגים שגויים, ייצוג שגוי של אנשים, ושימוש לא נכון בסוגי מזהים.
\par
התגובה הייתה לבנות מערכת מגינים מתועדת:
\begin{itemize}
  \item רשימה סגורה של מחלקות חוקיות עבור P31 על כתב יד (manuscript, codex, illuminated, composite, palimpsest), כך שלא ניתן להוסיף P31 כמו ``אדם'' או ``ארגון'' לפריט שזוהה ככתב יד.
  \item שמירה קפדנית של מאפייני זהות לפי קבוצות: שני אנשים עם זוגות זהות שונים (P227 GND, P214 VIAF, P569 תאריך לידה) לעולם לא ימוזגו.
  \item שלב אימות בארבעה שלבים בכל העלאה לוויקינתונים: רק פריטים שאני יוצרת אפשר לערוך, ובכל מקרה לפני כתיבה רעננים שוב.
  \item טיפול במזהים אנונימיים: שמות-מילוי כמו ``unknown'', ``לא ידוע'' או תארי תפקיד לא הופכים לפריטים, אלא נשמרים כטענת somevalue על הפריט של כתב היד.
  \item פסק זמן (moratorium) על העלאות המוניות לוויקינתונים, שניתן להסיר רק מתוך משתנה סביבה מפורש אחרי שכל הבדיקות עברו.
\end{itemize}
\par
כל מגן כזה גובה במבחני יחידה ובמבחני אינטגרציה אוטומטיים, נכון להיום למעלה מ-1{,}000 מבחנים שעוברים. ה-CLAUDE.md (קובץ ההנחיות הפנימי) מתעד את כל החוקים האלה ואת ההקשר ההיסטורי של כל אחד מהם.
\end{hebrew}


\section{\texthebrew{למה IIIF, ולא רק קישור ישיר לקובץ ה-RDF?}}
\begin{hebrew}
ה-RDF נשמר --- כל מַניפֶסט IIIF כולל מאפיין seeAlso שמוביל ל-RDF המלא. אבל IIIF הוא הכלי הנכון לסוג מסוים של מידע שאי אפשר לבטא במודל של ויקינתונים.
\par
IIIF מאפשר שלוש דברים שכרגע הם הפער המרכזי:
\begin{itemize}
  \item Canvas לכל פוליו: מבנה דף-אחר-דף שמאפשר להפנות לטענה לפי מיקום פיזי בכתב היד.
  \item Range לכל יחידה קודיקולוגית: ייצוג מובנה של מבנה הצירוף ולא רק כמספר חלקים.
  \item AnnotationCollection לקולופונים, התערבויות סופרים ושוליים: שכבות של הערות מקושרות למיקום בכתב היד.
\end{itemize}
\par
מודל הנתונים של ויקינתונים לא יכול לייצג את שלושת הדברים האלה בצורה טבעית --- אבל הוא יכול להחזיק קישור P6108 אחד לכל מַניפֶסט. מי שמתעניין בקריאה מעמיקה של הגרף יקבל את כל המבנה הזה דרך מציג IIIF סטנדרטי כמו Mirador, בלי לעזוב את עמוד הפריט בוויקינתונים.
\end{hebrew}


\section{\texthebrew{איך מפיקים תווית באנגלית לחיבור שכותרתו רק בעברית?}}
\begin{hebrew}
זאת שאלה לגיטימית כי ויקינתונים מצפה לתווית באנגלית לצורך גילוי גלובלי, אבל הרבה כתבי יד עבריים אין להם כותרת לטינית בקטלוג. תעתיק מכני נאיבי של עברית לא מנוקדת מפיק שלדים של עיצורים (למשל \texthebrew{תקנות רבנו גרשם מאור הגולה} הופך ל-``Tknot rvno grshm mor hgolh''), והקהילה של ויקינתונים דחתה את הצורה הזאת כלא ראויה לתוויות ציבוריות.
\par
התשובה היא מפל בן חמישה שלבים, שמעדיף מקורות סמכותיים על פני פלט אלגוריתמי:
\begin{itemize}
  \item שלב 1 --- מילון ידני קטן ($\sim$33 ערכים) של דמויות מפורסמות (רמב"ם, רש"י, רמב"ן, יוסף קארו, אבן עזרא) ומונחים נפוצים (תנ"ך, תלמוד, משנה).
  \item שלב 2 --- קריאת תעתיק ALA-LC ישירות מתוך שדות 246 ו-880 של רשומת ה-MARC. כאשר הספרנים של הספרייה הלאומית כבר ביצעו את התעתיק, אנחנו נשענים עליו.
  \item שלב 3 --- שאילתת SPARQL הפוכה: אם כבר קיים פריט בוויקינתונים שתווית העברית שלו זהה לקלט, אנחנו לוקחים את תווית האנגלית שלו. הצורה הקהילתית-המוסכמת ניצחת תמיד.
  \item שלב 4 --- מודל TaatikNet (\texttt{malper/taatiknet}), רשת ByT5-small שאומנה על כ-15 אלף זוגות עברית-לטינית מ-Wiktionary. המודל מפיק תעתיק שזורם פונטית (``Takanut rivno gereshem meor hagola'') ולא רק שלד עיצורים.
  \item שלב 5 --- טבלת ALA-LC עיצורית כתחתית אבסולוטית. למעשה מבוטלת בתוויות של חיבורים כי הפלט שלה מכוער מדי לפרסום ציבורי.
\end{itemize}
\par
הצורה הסופית של התווית היא מבנה משולב: \texttt{תעתיק (NLI מספר הרשומה)}, למשל ``Takanut rivno gereshem meor hagola (NLI 990001801390205171)''. החלק הראשון קריא; החלק השני מצביע בצורה חד-משמעית על רשומת המקור בקטלוג הלאומי. כאשר התעתיק נכשל לחלוטין (מודל לא זמין, קלט קצר מדי, או שהמודל החזיר תווים עבריים בטעות), התווית מצטמצמת ל-``NLI 990001801390205171'' בלבד --- מזהה יציב, חד-משמעי, שתמיד יוביל בלחיצה אחת לרשומת המקור.
\par
חוק קשיח: אסור לאף תווית באנגלית להכיל תווים בכתב עברי. הסיבה: ערבוב סקריפטים בתווית גלובלית מבלבל את משתמשי ויקינתונים ופוגע בחיפוש. אם TaatikNet החזיר תוצאה שמכילה תווים עבריים (תוצאה של ``הד'' של הקלט במקום תעתיק אמיתי), אנחנו דוחים את התוצאה כולה ונופלים אל מזהה ה-NLI לבדו.
\end{hebrew}


% =============================================================================
% Hard-tier questions — methodology, AI critiques, LLM comparisons
% =============================================================================


\section{\texthebrew{למה לא להשתמש פשוט ב-GPT-4 או Claude עם פלט מובנה במקום כל הפייפליין הזה?}}
\begin{hebrew}
זאת השאלה הראשונה שכל סטודנט שואל ב-2026, וצריך לענות עליה בכנות.
\par
ניסינו את זה. על מדגם קטן של עשר רשומות, GPT-4 עם הוראות מובנות הפיק תוצאות סבירות, אבל ארבע בעיות מהותיות מנעו אימוץ:
\begin{itemize}
  \item \textbf{עלות בקנה מידה}: 50 אלף רשומות × כ-15 אגורות לרשומה = כ-7,500 דולר לכל ריצה מלאה של הפייפליין. אצלנו זה אפס. המחקר חוזר על עצמו כשמשפרים מסווג, מוסיפים מקרה בוחן, או מאמתים שכבת סמכות חדשה.
  \item \textbf{שחזוריות מדעית}: המודל של OpenAI משתנה ללא הודעה. ערך F1 שמדדנו בינואר עלול להשתנות במרץ. דוקטורט שמסתמך על API חיצוני לא ניתן לשחזור בעוד חמש שנים. DictaBERT שמותקן לוקאלית — כן.
  \item \textbf{הזיות בעברית מובהקת}: על קולופונים בעברית רבנית עם קיצורים (``זצ"ל'', ``בכ"ר'', ``מהר"ם''), GPT-4 לעיתים המציא שמות שלא היו בטקסט. המודל שלנו, שאומן ישירות על קטלוג NLI, לא עושה זאת.
  \item \textbf{פרטיות וריבונות נתונים}: שליחת קטלוג הספרייה הלאומית של ישראל דרך API של חברה אמריקאית היא בעיה משפטית ואתית. הפייפליין שלנו רץ במלואו על מחשב נייד, ללא רשת.
\end{itemize}
\par
LLMs הם כלי משלים, לא תחליף. אנחנו משתמשים ב-TaatikNet (שהוא מודל קטן יותר), ובעתיד נשקול שילוב של מודל גדול לאימות, לא לחילוץ ראשוני.
\end{hebrew}


\section{\texthebrew{האם השוויתם פורמלית את ה-NER שלכם מול Gemini על אותו מדגם?}}
\begin{hebrew}
כן. בריצה החדשה השתמשנו במדגם ייצוגי של 100 רשומות, עם אותן תוויות זהב, והשווינו שלוש שיטות: המודל המאומן, Gemini בלי דוגמאות, ו-Gemini עם שלוש דוגמאות. התוצאות:
\begin{itemize}
  \item המודל המאומן: 88 true positives, 23 false positives, 20 false negatives, כלומר 80.4\% F1 בהתאמה קפדנית של שם ותפקיד.
  \item Gemini zero-shot: 67 true positives, 90 false positives, 41 false negatives, כלומר 50.6\% F1.
  \item Gemini three-shot: 68 true positives, 86 false positives, 40 false negatives, כלומר 51.9\% F1.
\end{itemize}
\par
החידוד החשוב הוא לא להתבלבל עם סוכן ההערכה. סוכן ההערכה שואל האם כל מועמד נראה סביר, ולכן הוא יכול לאשר הרבה מועמדים ש-Gemini מייצר. ה-F1 הקובע הוא מול תוויות הזהב: שם ותפקיד חייבים להתאים בדיוק, וחסרים נספרים כ-false negatives.
\end{hebrew}


\section{\texthebrew{DictaBERT מ-2022 — האם הוא לא מיושן ב-2026?}}
\begin{hebrew}
שאלה הוגנת. DictaBERT הוא אכן מודל מ-2022 שמבוסס על BERT, ועידני LLM קצרים. אבל היציבות שלו היא יתרון, לא חיסרון:
\begin{itemize}
  \item הוא רץ במלואו על מחשב נייד M4, ב-MPS. ללא רשת, ללא תלות בעלות.
  \item הוא לא משתנה. הציון שאני מדווח עליו היום יהיה זהה בעוד שנתיים.
  \item ארכיטקטורת BERT לתיוג רצפים (sequence labeling) היא עדיין SOTA למשימות NER ספציפיות-לדומיין כאשר יש דאטה מאומן. ML papers בכנס ACL 2025 ממשיכים להציג אדריכלויות מבוססות-BERT לעברית עם ביצועים תחרותיים.
  \item DictaBERT אומן ספציפית על עברית. מודלים גדולים יותר מבוססי decoder (Falcon, Llama) פחות מדויקים על תיוג רצפי-מילה ספציפי לעברית.
\end{itemize}
\par
גם בדקנו את זה אמפירית מול מודל חדש יותר: \textenglish{NeoDictaBERT bilingual}. על אותה משימת \textenglish{Provenance NER}, עם אותו מערך של 12,100 דוגמאות וחלוקת חמשת הקיפולים, \textenglish{DictaBERT} הגיע ל-\textenglish{95.91\% F1} בממוצע, בעוד \textenglish{NeoDictaBERT} הגיע ל-\textenglish{94.65\% F1}. ההשוואה לא מושלמת, כי הריצה החדשה השתמשה ב-batch קטן יותר ובתקציב אימון קצר יותר כדי לשמור על יציבות ב-\textenglish{MPS}, אבל היא כן מחזקת את הטענה שלא מחליפים מודל רק בגלל שהוא חדש יותר.
\par
אם DictaBERT-3 ישוחרר ויעבוד טוב יותר, נחליף בקלות. הפייפליין שלנו מצביע על הקובץ הפיזי של המודל; שדרוג הוא שינוי שורה אחת.
\end{hebrew}


\section{\texthebrew{ארבעה מודלים קטנים — האם זה לא לא יעיל בהשוואה ל-LLM אחד מותאם?}}
\begin{hebrew}
שתי תשובות:
\par
ראשית, יעילות תפעולית. כל מודל קטן רץ עצמאית, אפשר לאמן ולעדכן כל אחד בלי לגעת באחרים. אם שיפרתי את מודל הבעלות אין צורך לעבור על מבחני האנשים. ב-LLM יחיד, כל שינוי נוגע בכל הפלטים.
\par
שנית, פירוש (interpretability). כאשר טעות מתרחשת, יש לי ארבע נקודות חיתוך מובנות --- NER לאנשים, NER לבעלות, NER לתוכן, ומסווג הז'אנר. אני יודע אם הטעות הייתה במסווג ז'אנר, או בחיבור לסמכות, או באחד ממודלי ה-NER. ב-LLM אחד אני לא יודע. הפרובננס שאני מדבר עליה במצגת --- שכל עובדה יודעת מאיפה היא הגיעה --- מתאפשרת רק במערכת רכיבים מובחנת.
\par
לבסוף, יציבות: כשמודל אחד נכשל (למשל TaatikNet מחזיר עברית במקום לטינית), הפייפליין ממשיך עם פולבק. ב-LLM יחיד אין פולבק טבעי.
\par
אז כן, מבחינת זמן GPU בלבד, LLM אחד אולי יותר יעיל. אבל מבחינת תחזוקה, ביקורת, ופירוש --- ארבעה מודלים ייעודיים הם הבחירה הנכונה למחקר אקדמי.
\end{hebrew}


\section{\texthebrew{למה בכלל להתאמן? GPT-4 עם הוראות סכמה ברורות כנראה ישיג את ה-F1 שלכם בלי אימון.}}
\begin{hebrew}
על משימות סטנדרטיות אנגליות, התשובה הזאת קרובה לאמת. על עברית רבנית, היא לא נכונה.
\par
הקטלוג של ה-NLI מכיל קונבנציות שגיבשו ספרנים במשך 75 שנה: ``בכ"ר'' כסימן ל-בן כבוד הרב, ``מהר"ם'' כראשי תיבות לרב מסוים, ``זצ"ל'' כסימן לאדם שנפטר, ``המעתיק'' כתפקיד מובהק שמובחן מ-``המגיה''. GPT-4 מכיר חלק מהסימנים האלה אבל לא את כולם. ב-zero-shot הוא לא יודע לתת לרב יוסף בן מאיר את התפקיד הנכון לעומת הסופר חיים בן יוסף בקולופון.
\par
המודל שלנו אומן ספציפית על הקטלוג הזה, על הקיצורים האלה, על הסגנון הזה. במבחן החדש, מול תוויות זהב ולא מול שופט LLM, המודל המאומן הגיע ל-80.4\% F1 בהתאמה קפדנית של שם ותפקיד, 86.8\% F1 בחילוץ שם בלבד, ו-92.6\% דיוק תפקיד כאשר השם זוהה. באותו מדגם, Gemini zero-shot הגיע רק ל-50.6\% F1 קפדני, ו-Gemini עם שלוש דוגמאות הגיע ל-51.9\%.
\par
ואם אנחנו רוצים F1 גבוה יותר עם LLM? צריך לתת לו כמה דוגמאות (few-shot). וכמה דוגמאות זה כבר אימון, רק ביוקר.
\end{hebrew}


\section{\texthebrew{מה עלות הפליטות (carbon) של ההכנה החלוצה והאימון בקנה מידה זה?}}
\begin{hebrew}
שאלה חשובה ולא תמיד נשאלת. אענה בכנות:
\begin{itemize}
  \item אימון כל מודל NER: כ-3-4 שעות על M4 Pro = כ-0.05 קוט"ש = כ-25 גרם CO2 (ישראל, רשת חשמל מעורבת).
  \item אימון של כל החמישה: כ-15 שעות סך הכל = כ-125 גרם CO2.
  \item הפליטות הזעירות האלה מתבטלות מול הפלט: רישומי מחקר שמשתמשים בהם חוקרים בעוד 20 שנה.
  \item השוואה: ריצה אחת של GPT-4 על 50 אלף רשומות מצרכת כ-450 גרם CO2 (אומדן OpenAI). כל ריצה חוזרת מוסיפה.
\end{itemize}
\par
המחקר שלנו זול אקולוגית בגלל שתי החלטות מודלי-מודלים: מודלים קטנים שרצים על מחשב נייד, ושמירה של הנתונים בלוקאלית. הברירה האלטרנטיבית — LLM גדול ב-API — היא דווקא בעלת עלות סביבתית גבוהה יותר בקנה מידה.
\end{hebrew}


\section{\texthebrew{68 רשומות זה לא מדגם משמעותי סטטיסטית. איך אתם מגינים על כל טענה כמותית מתוך מדגם כל כך קטן?}}
\begin{hebrew}
זה הביקורת הקשה ביותר, ואני מודה בה בפתיחות.
\par
ראשית, 68 רשומות הן לא מדגם סטטיסטי במובן של null-hypothesis testing. הן \textbf{מדגם איכותי לעיון לעומק}: כל רשומה נבדקת ידנית, וכל טענה כמותית קשורה לרשומות ספציפיות שאפשר לבחון.
\par
שנית, 68 רשומות הן מספיק בגלל הטרוגניות: לא בחרנו את הקלות. במדגם יש כתבי יד שלמים וחלקיים, עברית וארמית, רשומות עם NER הצלחה מלאה ורשומות עם כשלון, רשומות נדירות כמו כתב יד אחד שמופיע במצב מיוחד.
\par
שלישית, ההצלחה היחסית כן ניתנת להעברה לקורפוס המלא. כיסוי תאריכים של 97\% אינו תיאוריה — זה תוצאה של חוקים שאפשר להריץ על כל הקורפוס ולמדוד.
\par
בעתיד? כן, אני מסכים שיוכל להיות נכון להרחיב את מדגם ההערכה ל-300 או 500 רשומות. זה משימה שדורשת מאה שעות סקירה ידנית, ואני מקוה לעשות זאת אחרי הדוקטורט עם תרומה של חוקרים נוספים.
\end{hebrew}


\section{\texthebrew{השגחה מרחוק מייצרת תוויות רועשות. מדדתם את שיעור הרעש?}}
\begin{hebrew}
כן. הרעש בדאטה הוא הבעיה המרכזית של distant supervision, ויש לנו שלוש שכבות התמודדות:
\begin{itemize}
  \item \textbf{מדידה ישירה}: על מדגם אקראי של 100 דוגמאות אימון אנושית, מדדנו רעש של כ-12-15\% עבור מודל הבעלות, וכ-8\% עבור מודל האנשים. המדידה נמצאת בקובץ \texttt{ner/dictabert\_error\_analysis.json}.
  \item \textbf{פילטרים פוסט-NER}: אחרי המודל יש שכבת \texttt{converter/authority/ner\_post\_filters.py} שמסירה false-positives ידועות --- שמות מקומות שזוהו כאנשים, תאריכי דפים שזוהו כתאריכים אמיתיים, וכן הלאה.
  \item \textbf{סקירת אדם}: עומס הרעש שמגיע לסוף הוא ~3-5\%, ואת זה האדם מסיר בממשק.
\end{itemize}
\par
האם 12\% רעש באימון מסכן את המודל? לא, ובדקנו זאת. המודל לומד לפעמים מדוגמה שגויה, אבל הוא לומד הרבה יותר מ-88\% הדוגמאות הנכונות. F1 של 80.4\% (קפדני, שם+תפקיד) ו-86.8\% (שם בלבד) על מערך הזהב מראה שזה עובד.
\end{hebrew}


\section{\texthebrew{מודל התוכן ב-99.99\% F1 — זה דגל אדום של התאמת יתר (overfitting). מה קורה שם?}}
\begin{hebrew}
ביקורת מצוינת ואני מסכים שזה צריך הסבר.
\par
הציון הזה אמיתי, אבל הוא מטעה. שתי סיבות:
\begin{itemize}
  \item \textbf{מדגם הערכה זעיר}: רק 2 מתוך 68 רשומות מכילות שדה MARC 505 (תוכן). הציון מחושב על אותן שתי רשומות. F1 גבוה על מדגם של 2 הוא מספר אבל לא ראיה.
  \item \textbf{משימה מובנית במיוחד}: שדה 505 הוא חצי-מובנה מלכתחילה. הפורמט הוא ``דף 1א-50ב: שם החיבור, מאת מחבר''. המודל למד תבנית פשוטה ויציבה, לא חופש לשוני אמיתי.
\end{itemize}
\par
לכן בפרק התוצאות אני מתאר את הציון הזה לא כ-99\% F1 מודלי, אלא כ-``מודל שלמד תבנית פורמלית ומעתיק אותה היטב''. אם מחר תגיע רשומה עם שדה 505 בפורמט חדש שלא ראינו, המודל ייכשל ביעילות. זאת מגבלה של distant supervision שאני מודה בה.
\par
המודל המעניין יותר הוא מודל האנשים, שעובד על טקסט חופשי לגמרי. במבחן החדש הוא מגיע ל-80.4\% F1 בהתאמת שם ותפקיד יחד, ול-86.8\% F1 כאשר מודדים חילוץ שם בלבד. שם הציון משקף יכולת אמיתית, כי הוא נמדד מול תוויות זהב ולא מול שיפוט סובייקטיבי של LLM.
\end{hebrew}


\section{\texthebrew{האם ה-eval set שלכם דלוף לדאטה של האימון דרך התאמת distant supervision?}}
\begin{hebrew}
שאלה מצוינת וביקורת מתודולוגית חשובה. התשובה היא: כמעט בטוח שיש זליגה חלקית, ואני מודה בה.
\par
מנגנון distant supervision שואב דוגמאות אימון מקטלוג ה-NLI. כל הקורפוס נמצא בקטלוג. ה-68 רשומות הן תת-קבוצה של הקטלוג. אם רשומה ב-68 מכילה את השם ``משה בן מימון'' בהערה חופשית, וגם בשדה מובנה, היא תהפוך לדוגמת אימון.
\par
כיצד התמודדנו? בנינו הפרדה: לאחר חילוץ דוגמאות distant supervision, השווינו את control numbers של 68 הרשומות מול דאטה האימון. כל דוגמה ש control number שלה שייך ל-68 — הוסרה. זה ב-\texttt{scripts/extract\_genre\_samples.py}.
\par
האם זה מספיק? לא לחלוטין. ייתכן שדוגמה אחרת בקורפוס מכילה משפט דומה. אבל הסיכון של זליגה ספציפית-לרשומה — הוא נמוך.
\par
התיקון הסופי: בעתיד אני מתכוון לבנות eval set מ-200 רשומות שמעולם לא נכנסו לקטלוג ה-NLI עצמו — למשל כתבי יד שטרם נוקטלגו במלואם. זה ייבטל את הזליגה הזאת לגמרי.
\end{hebrew}


\section{\texthebrew{איפה לימודי הפילוח (ablation studies)? מה התרומה השולית של כל רכיב?}}
\begin{hebrew}
שאלת מתודולוגיה חזקה. עשינו פילוח חלקי, אבל לא מקיף, ואני מודה בכך:
\begin{itemize}
  \item \textbf{פילוח שעשינו}: השוואת כיסוי לפני ואחרי הוספת חוקים (תאריכים: 22\% → 97\%); לפני ואחרי NER (שמות אנשים: 0 → 59 מופעים); לפני ואחרי המסווג (ז'אנרים: 28 → 39).
  \item \textbf{פילוח שלא עשינו}: השפעה של מספר השכבות של DictaBERT שמוקפאות, השפעה של גודל הקלט (max\_length=64 לעומת 128), השפעה של temperature ב-TaatikNet, השפעה של ספר הסמכות (Mazal לבד, Mazal+VIAF, וכו').
\end{itemize}
\par
מדוע לא? כי הפילוחים האלה דורשים עוד אימונים חוזרים — כל אחד 3-4 שעות. בתקופת הדוקטורט בחרנו להשקיע את הזמן הזה בשיפור הקטלוג, לא בלימודי פילוח. זאת מגבלה שאני מציין במפורש בפרק המתודולוגיה.
\par
לפרסום בכתב עת אקדמי (SWJ), אני מתכוון להוסיף לפחות שלוש סטים של פילוח עיקריים: השפעת מספר השכבות הקפואות, השפעת גודל מאגר הסמכות, והשפעת אורך הקלט.
\end{hebrew}


\section{\texthebrew{למה פייפליין של 6 שלבים ולא end-to-end joint training?}}
\begin{hebrew}
שאלה ארכיטקטונית טובה. שלוש סיבות:
\par
\textbf{ראשית, פירוש}. כל שלב מייצר פלט שאני יכול לבדוק עצמאית. אם כיסוי הסמכויות נמוך, אני יודע אם הבעיה ב-NER (חילוץ שמות שגויים) או בשלב הסמכויות (התאמה שגויה). end-to-end model הוא קופסה שחורה.
\par
\textbf{שנית, חזרה שולית}. אם השתפר מודל NER, אני לא צריך לחזור על שלב הסמכויות מאפס. אם המסווג ז'אנר עודכן, רק ספציפית הפלט שלו משתנה. ב-end-to-end, כל שינוי דורש אימון מלא.
\par
\textbf{שלישית, מקורות שונים, שכבות שונות}. NER מתעניין בטקסט; הסמכויות מתעניינות בזיהוי; ה-RDF מתעניין בסמנטיקה. אלה תחומים שונים שמתאימים לאדריכלויות שונות. בעבר אכן ניסו end-to-end לפייפליינים דומים (למשל DBpedia Spotlight) והגיעו למסקנה שזה לא עובד טוב על דאטה הטרוגני.
\par
לבסוף, הוגנות: אכן end-to-end model אקדמי תיאורטי יכול להיות יפה. אבל הוא לא מציאותי לעבודה ספרנית מעשית עם תהליך סקירה ידני.
\end{hebrew}


\section{\texthebrew{למה אפליקציית Qt למחשב במקום שירות web?}}
\begin{hebrew}
ארבע סיבות מהותיות:
\par
\textbf{עבודה ללא רשת}. ספרנים בארכיבים עובדים לעיתים במחסנים, באולמות קריאה, או במקומות עם אינטרנט מוגבל. אפליקציה מקומית עובדת בכל מקום.
\par
\textbf{ביצועים}. הפייפליין רץ על מודלים גדולים (~5 גיגה של נתוני מודל). שליחת הכל לשרת — או לקוח-בענן — כרוך בזמן רשת מאסיבי. מקומית זה מיידי.
\par
\textbf{פרטיות}. הקטלוג של NLI נמצא במחשב המקומי. אנחנו לא רוצים לשלוח אותו לשרת חיצוני.
\par
\textbf{עצמאות ארוכת טווח}. שירות web דורש תחזוקת שרת, אישורים, ועלויות תפעוליות. אפליקציה ניידת רצה לעד, בלי תלות בתקציב.
\par
אם בעתיד נרצה שיתוף — אנחנו מתכוונים לפרסם את הפלט כקובץ סטנדרטי (RDF/JSON/QuickStatements) שמכל אחד יכול להעלות לכלים שלו. הפייפליין הוא הכלי המקומי; הפלט הוא ה-LOD המשותף.
\end{hebrew}


\section{\texthebrew{human-in-the-loop מגביל קנה מידה. איך זה יסתרג מעבר לדוקטורט?}}
\begin{hebrew}
ביקורת חכמה ואני מסכים. ל-human-in-the-loop יש שתי קומות שונות:
\par
\textbf{הקומה הראשונה היא הסקירה הראשונית של כל ההפלט החדש}. זה אכן צוואר בקבוק. עבור 68 רשומות בדוקטורט שלי, אני סקרתי ידנית הכל. עבור 50 אלף רשומות, זה לא אפשרי.
\par
\textbf{הקומה השנייה היא סקירת המקרים הגבוליים בלבד}. כאן הפייפליין מסייע: כל טענה מקבלת ציון ביטחון. טענות עם ביטחון גבוה (95\% +) מאושרות אוטומטית; טענות עם ביטחון בינוני (60-95\%) דורשות אישור; טענות עם ביטחון נמוך נדחות.
\par
המתודולוגיה הזאת מאפשרת לחוקר אחד לאשר 1000 רשומות ביום במקום 50. עדיין לא קנה מידה של מיליונים, אבל סדר גודל יותר ממה שהקטלוג היה מסוגל לקלוט.
\par
לטווח ארוך, השפעה אמיתית באה משילוב של (1) מודלים שנעשים בטוחים יותר עם משוב, (2) חוקרי ספרייה שמיומנים בעבודה עם הפייפליין, ו-(3) פיתוח של ציוני ביטחון יותר טובים. הדוקטורט הוא הצעד הראשון, לא הפתרון הסופי.
\end{hebrew}


\section{\texthebrew{HMO + ויקינתונים + ויקיבייס — שלושה פלטים זה שלושה עומסי תחזוקה.}}
\begin{hebrew}
מסכים שזה עומס משמעותי. הסיבה לבחירה הזאת היא בכוונה תחילה:
\par
\textbf{HMO RDF} הוא הפלט המחקרי המלא. הוא נשמר כקובץ TTL סטטי שלא דורש תחזוקה — פעם שהוא נוצר, הוא קיים. אין שרת לתחזק.
\par
\textbf{הוויקיבייס של הפרויקט} (mhm-hmo.wikibase.cloud) הוא ממשק חיפוש לצפייה אנושית. הוא נוצר על ידי Wikibase Cloud, שירות חינמי של Wikimedia Foundation, ולא דורש תחזוקת שרת.
\par
\textbf{ויקינתונים} הוא הקצה הציבורי הרחב. ההקרנה הזאת היא חד-כיוונית: אנחנו דוחפים נתונים, ולא מתחזקים שום שירות.
\par
אז למרות שיש שלוש שכבות, הכבדה התפעולית האמיתית היא לא שלוש פעמים — היא בערך פעם וחצי. הקובץ TTL מטופל אוטומטית, Wikibase Cloud מטופל על ידי Wikimedia, ורק העדכון של ויקינתונים דורש זמן אדם.
\par
המחיר הזה הגיוני: יש לנו פלט מחקרי עמוק (HMO), פלט נווט (ויקיבייס), ופלט ציבורי (ויקינתונים) שכל אחד משרת קהל יעד שונה.
\end{hebrew}


\section{\texthebrew{מה החדשנות האמיתית פה מול OpenRefine + ממפה CIDOC?}}
\begin{hebrew}
זאת השאלה הקשה ביותר ועמדתי בפניה הרבה. תשובתי הוגנת:
\par
OpenRefine + ממפה CIDOC יכולים לעשות חלק מהעבודה: התאמת מחרוזות לסמכויות, מיפוי של שדות MARC לתכונות CIDOC, ניקוי טבלאות. אבל הם לא יכולים לעשות שלושה דברים שהפייפליין שלנו עושה:
\begin{itemize}
  \item \textbf{NER על טקסט חופשי בעברית}. אין מסלול ב-OpenRefine לחילוץ ``שלמה בן יצחק'' מתוך משפט ``בעלים: שלמה בן יצחק זצ"ל''. צריך מודל אומן.
  \item \textbf{מבנה HMO קודיקולוגי}. CIDOC-CRM הוא כללי. HMO מוסיף שכבה ייעודית — יחידות קודיקולוגיות, התערבויות סופרים, קולופונים ברמת פוליו — שאין מקבילה ישירה ב-CIDOC.
  \item \textbf{הקרנה זהירה לויקינתונים}. אחרי האירוע מול הקהילה באפריל 2026, אנחנו מבינים שהקרנה לויקינתונים מצריכה שכבת אימות בארבעה שלבים. OpenRefine לא מציע זאת.
\end{itemize}
\par
החדשנות שלנו אינה בכל רכיב יחיד, אלא ב\textbf{שילוב}: NER הספציפי לעברית רבנית + אונטולוגיית HMO + שכבת הגשר לויקינתונים. כל אחד מהשלושה כבר קיים בנפרד; הסינתזה שלהם — לא.
\par
שנית, פלט מעשי: הקטלוג של NLI כיום מקבל את הפלט שלנו, וזה הצעד הראשון לעיבוד מאסיבי של הקורפוס שלא קיים אצל אף כלי קודם.
\end{hebrew}

\end{document}
